\documentclass[a4paper,12pt,fleqn]{article}%fleqn pour aligner les equations à gauche

\usepackage[french]{babel}
\usepackage[sf,tiny,bf]{titlesec}
\usepackage[normalem]{ulem}
\usepackage{romannum}
%\allsectionsfont{\sffamily\raggedright\underline}






\usepackage{tgheros,textcomp}%Police Arial
\renewcommand{\familydefault}{\sfdefault}

%\setlength{\mathindent}{0pt}%indentation dans math = 0

\usepackage[top = 1cm, bottom = 2cm, left = 1cm, right = 1cm]{geometry}
%\usepackage{titlesec}
%\setcounter{secnumdepth}{4}
\usepackage{xcolor}
\usepackage{amssymb} %double flèche de réaction
%\usepackage{xtab}%tableau sur plusieurs pages
%\usepackage{esvect}%grands vecteurs

\usepackage[version=4]{mhchem} %equations chimiques

\usepackage{array}
\usepackage{chemfig}
\usepackage{multirow}
%\usepackage{sectsty}
%\usepackage{caption}
%\usepackage{amsmath}
%\usepackage{amssymb}
\usepackage[cal=boondoxo]{mathalfa} 
\usepackage{enumitem}
\usepackage{lastpage}%Pour noter le bon nombre de pages au total
\usepackage{tcolorbox}
\usepackage{siunitx}
\usepackage[european, straightvoltages, RPvoltages]{circuitikz}
\usepackage{tikz}
\usetikzlibrary{shapes.geometric}
\usetikzlibrary{calc}
\usetikzlibrary{automata,positioning}
\usepackage{multicol}
\usepackage{eqnarray}
\usepackage{tabularx,colortbl}%colorer fond cellules

\sisetup{
	detect-all,
	output-decimal-marker={,},
	group-minimum-digits = 3,
	group-separator={~},
	number-unit-separator={~},
	inter-unit-product={~}
} %setup pour SIunitx


\usepackage{pgf}
\usepgflibrary{fpu}
\pgfkeys{/pgf/number format/.cd, precision = 2, fixed, fixed zerofill}

\usepackage{listings}%pour insérer le code python
\definecolor{codegreen}{rgb}{0,0.6,0}
\definecolor{codegray}{rgb}{0.5,0.5,0.5}
\definecolor{codepurple}{rgb}{0.58,0,0.82}
\definecolor{backcolour}{rgb}{0.95,0.95,0.92}


\titlespacing\section{5pt}{5pt plus 4pt minus 2pt}{5pt plus 2pt minus 2pt}
\titlespacing\subsection{8pt}{5pt plus 4pt minus 2pt}{5pt plus 2pt minus 2pt}

%\sectionfont{\fontsize{14}{8}\selectfont}
\renewcommand{\thesection}{\fontsize{14}{8}\arabic{section}. }
\renewcommand{\thesubsection}{\hspace{2cm}\emph{\alph{subsection}.}}

\title{\vspace{-1cm}\large 
	\begin{tabular}{|>{\centering}m{6cm}|>{\centering}m{7cm}|>{\centering\arraybackslash}m{4cm}|}
		\hline
		Questionnaire & \textbf{\textsc{Préparation à la mécanique}} &  Chapitre 4 \\
		\hline
\end{tabular}}
\date{}

\newpagestyle{mystyle}{%
	\setfoot{ROSALIE}{\thepage\,$|$\,\pageref{LastPage}}{}
}
\renewpagestyle{plain}{%
	\setfoot{ROSALIE}{\thepage\,$|$\,\pageref{LastPage}}{}
}

\pagestyle{mystyle}

\newenvironment{itemize*}%
{\begin{itemize}[label=$-$]%
		\setlength{\itemsep}{0pt}%
		\setlength{\parskip}{0pt}}%
	{\end{itemize}}

\newenvironment{itemize**}%
{\begin{itemize}[label=]%
		\setlength{\itemsep}{0pt}%
		\setlength{\parskip}{0pt}}%
	{\end{itemize}}

\begin{document}
	\maketitle
	\vspace{-2.2cm}
\begin{enumerate}
	\item Quels sont les éléments nécessaires à la description d'un mouvement ?
	\item Définir la vitesse ?
	\item Etablir les coordonnées des vecteurs suivants :
\end{enumerate}
\begin{tikzpicture}
	\draw[thick, -Stealth] (0,0) to (0,5.5);
	\draw[thick, -Stealth] (0,0) to (5.5,0);
	\foreach \k in {0,1,...,5} {
		\draw[thick] (\k,-0.1 ) -- (\k,0.1);
		\draw[thick] (-.1,\k) -- (0.1,\k);
		}
	\node[left] (y) at (0,5){$y$};
	\node[below] (x) at (5,0){$x$};
	
	\draw[thick,blue!60!green,-Stealth] (2,1) to (2,4);
\end{tikzpicture}
\begin{tikzpicture}
	\draw[thick, -Stealth] (0,0) to (0,5.5);
	\draw[thick, -Stealth] (0,0) to (5.5,0);
	\foreach \k in {0,1,...,5} {
		\draw[thick] (\k,-0.1 ) -- (\k,0.1);
		\draw[thick] (-.1,\k) -- (0.1,\k);
		}
	\node[left] (y) at (0,5){$y$};
	\node[below] (x) at (5,0){$x$};
	
	\draw[thick,blue!60!green,-Stealth] (0,0) to (2,3);
\end{tikzpicture} \\↑

\begin{tikzpicture}
	\draw[thick, -Stealth] (0,0) to (0,5.5);
	\draw[thick, -Stealth] (0,0) to (5.5,0);
	\foreach \k in {0,1,...,5} {
		\draw[thick] (\k,-0.1 ) -- (\k,0.1);
		\draw[thick] (-.1,\k) -- (0.1,\k);
		}
	\node[left] (y) at (0,5){$y$};
	\node[below] (x) at (5,0){$x$};
	
	\draw[thick,blue!60!green,-Stealth] (1,3) to (3,3);
\end{tikzpicture}
\begin{tikzpicture}
	\draw[thick, -Stealth] (0,0) to (0,5.5);
	\draw[thick, -Stealth] (0,0) to (5.5,0);
	\foreach \k in {0,1,...,5} {
		\draw[thick] (\k,-0.1 ) -- (\k,0.1);
		\draw[thick] (-.1,\k) -- (0.1,\k);
		}
	\node[left] (y) at (0,5){$y$};
	\node[below] (x) at (5,0){$x$};
	
	\draw[thick,blue!60!green,-Stealth] (1,3) to (4,1);
\end{tikzpicture}
\begin{enumerate}[resume]
	\item Etablir la somme des vecteurs suivants :\\
	\begin{tikzpicture}
	\draw[thick, -Stealth] (0,0) to (0,5.5);
	\draw[thick, -Stealth] (0,0) to (5.5,0);
	\foreach \k in {0,1,...,5} {
		\draw[thick] (\k,-0.1 ) -- (\k,0.1);
		\draw[thick] (-.1,\k) -- (0.1,\k);
		}
	\node[left] (y) at (0,5){$y$};
	\node[below] (x) at (5,0){$x$};
	
	\draw[thick,blue!60!green,-Stealth] (1,3) to (4,1);
	\draw[thick,red!60!orange,-Stealth] (2,4) to (4,4);
\end{tikzpicture}
\begin{tikzpicture}
	\draw[thick, -Stealth] (0,0) to (0,5.5);
	\draw[thick, -Stealth] (0,0) to (5.5,0);
	\foreach \k in {0,1,...,5} {
		\draw[thick] (\k,-0.1 ) -- (\k,0.1);
		\draw[thick] (-.1,\k) -- (0.1,\k);
		}
	\node[left] (y) at (0,5){$y$};
	\node[below] (x) at (5,0){$x$};
	
	\draw[thick,blue!60!green,-Stealth] (1,1) to (3,1);
	\draw[thick,red!60!orange,-Stealth] (2,2) to (2,5);
\end{tikzpicture}
	\item En se basant  sur les trajectoires suivantes, représenter le vecteur vitesse aux points marqués par une croix :
\end{enumerate}

\begin{tikzpicture}
	\foreach \k in {0,1,...,5} {
		\draw[thick] (\k,0.1 ) -- (\k,0.1);
		\draw[thick] (-.1,\k) -- (0.1,\k);
		}
	\node[left] (y) at (0,5){$y$};
	\node[below] (x) at (5,0){$x$};
\end{tikzpicture}
\end{document}